\section{MCP：外部工具集成}
%% ============================================================

\subsection{什么是 MCP}

Model Context Protocol（MCP）是一个开放标准，用于将 Codex 连接到外部工具和系统。当 Codex 所需的上下文不在代码仓库中时，MCP 就是解决方案。

\begin{importantbox}{何时使用 MCP}
\begin{itemize}[nosep]
    \item 所需上下文在仓库之外（如 issue tracker、monitoring 系统）
    \item 数据频繁变化（如实时指标、日志）
    \item 希望 Codex 使用工具而非依赖粘贴的指令
    \item 需要跨用户、跨项目的可复用集成
\end{itemize}
\end{importantbox}

\subsection{MCP 的配置方式}

Codex 支持两种 MCP server 类型：

\begin{itemize}
    \item \textbf{STDIO}：标准输入/输出通信
    \item \textbf{Streamable HTTP}：支持 OAuth 认证的 HTTP 通信
\end{itemize}

配置入口：
\begin{itemize}
    \item Codex App：Settings $\rightarrow$ MCP servers
    \item CLI：使用 \texttt{codex mcp add} 命令添加自定义 server
\end{itemize}

\begin{warningbox}{不要一次性接入所有工具}
只在工具确实能解锁某个真实工作流时才添加。不要一开始就把所有常用工具都接入。从 1--2 个明确能消除手动重复操作的工具开始，然后逐步扩展。
\end{warningbox}

\subsection{本章小结}

MCP 是 Codex 连接外部系统的标准协议，支持 STDIO 和 Streamable HTTP 两种模式。按需接入、从小处开始是正确的使用策略。

%% ============================================================
